\textbf{(OT01.1)} The four stars are circled and labelled on the sky map:
S3 ($\beta$ Peg, $23^h3^m47^s$, $28^\circ4'58''$) at top, S1 ($\alpha$ And,
$0^h8^m24^s$, $29^\circ5'16''$) to its lower left on the map, S2
($\alpha$ Peg, $23^h4^m46^s$, $15^\circ12'17''$) at the centre and S4
($\gamma$ Peg, $0^h13^m14^s$, $15^\circ10'59''$) below.

% FIGURE NEEDED: marked sky map with S1-S4 circled and labelled and the
% position of Z constructed and marked with the tolerance rectangles
% (2025_Map-OT01-Solution.pdf, single-page overlay; also reproduced on
% 2025_TEL_sol.pdf pages 1-2)

\begin{center}
\begin{tabular}{l r}
\hline
Circle around each correct star & 0.5 \\
Correct identification as S1--S4 of each star & 1.0 \\
\hline
\end{tabular}
\end{center}

\textbf{(OT01.2)}
\begin{itemize}
\item Realize that two stars have almost the same RA and two of them have
  almost same declination.
\item We draw a line through them indicating constant RA and constant
  Declination. (On the official overlay the construction lines are labelled
  $\delta \approx 15^\circ$, $\alpha \approx 23^h$, and their parallels
  through Z are labelled $\alpha \approx 21^h36^m$ and
  $\delta \approx -26^\circ$.)
\item Given the scale of the map and the separation between the 4 stars and
  the new object we can assume that the celestial grid will be almost
  linear.
\item We can thus mark `Z' at its given coordinates
  ($\alpha = 21^h36^m$; $\delta = -26^\circ10'$).
\end{itemize}

\begin{center}
\begin{tabular}{l r}
\hline
Object marked within inner rectangle & 7.0 \\
\quad Object marked between outer and inner rectangles & 4.0 \\
\hline
\end{tabular}
\end{center}

\textbf{(OT01.3)}
\begin{itemize}
\item After marking the new object on the given sky map, we carefully look
  at its position and the star pattern surrounding it.
\item Using the 25 mm focal length eyepiece we try to locate it.
\item By doing some fine adjustments, we try to bring the same pattern in
  the 10 mm focal length eye piece (which has a crosshair).
\end{itemize}

\begin{center}
\begin{tabular}{l r}
\hline
Pointing towards the screen & 1.0 \\
\hline
\multicolumn{2}{l}{If 10 mm objective is used:} \\
% sic: the official marking table says "objective" for the eyepiece
\quad Correct choice of eyepiece & 2.0 \\
\quad Object is anywhere in FOV & 3.0 \\
\quad Any part of the object is touching the origin of crosshair & 6.0 \\
\hline
\multicolumn{2}{l}{If 25 mm objective is used:} \\
\quad Object is anywhere in FOV & 3.0 \\
\hline
\end{tabular}
\end{center}
